\section{Datos y supervisión débil}
\label{sec:datos}

\textbf{Conjunto y reducción de escala.} El conjunto tiene 4\,407 estudios, 24\,371 series y
819\,078 archivos DICOM (265\,GB comprimidos; 570\,GB descomprimidos). Los scripts leen los
DICOM directamente desde el zip sin descomprimirlo y extraen solo las series necesarias,
unos 11\,GB.
\begin{itemize}
  \item \textbf{Series:} por estudio se elige una serie sagital con supresión de grasa,
        sensible a líquido. Se prefieren series 2D de 15 a 60 cortes; una adquisición 3D solo
        se usa si no hay otra, y se submuestrea a 32 cortes. La serie coronal con supresión
        de grasa se reserva para la ablación de plano.
  \item \textbf{Muestra:} 600 estudios no dorados, elegidos con cupos de al menos 60
        positivos de alta confianza por hallazgo y de al menos 15 estudios por idioma, y
        completados al azar. Se dividen en \texttt{weak\_train} (510) y \texttt{weak\_val}
        (90) eligiendo, entre 300 particiones aleatorias, la de menor diferencia de
        prevalencia (como máximo 3,9 puntos).
  \item \textbf{Dorados:} los 58 forman \texttt{gold\_val} y solo se usan en la evaluación
        final.
\end{itemize}

\textbf{Preprocesamiento.} Los cortes se ordenan por su posición a lo largo de la normal del
plano. La intensidad se normaliza por volumen entre los percentiles 0,5 y 99,5, se rellena a
cuadrado y se escala a 256$\times$256 (recorte de 224 al entrenar). Se conservan todos los
cortes, con una mediana de 29: el menisco medial y el lateral están en extremos opuestos de
la pila sagital, así que quedarse solo con los cortes centrales los perdería. Dos dorados sin
serie con supresión de grasa usan la sagital sin supresión; los resultados se reportan también
sin ellos.

\textbf{Extractor de etiquetas.} Detecta el idioma, normaliza el texto (Unicode, mayúsculas,
diacríticos y saltos de línea forzados) y elimina la historia clínica. Luego separa el cuerpo
de la conclusión y busca términos por hallazgo en nueve idiomas, siguiendo las definiciones
oficiales de cada etiqueta: por ejemplo, ``menisco'' significa rotura y ``contusión'' significa
edema óseo postraumático. La negación se detecta al estilo NegEx~\citep{chapman2001negex}:
siete palabras antes del término o cinco después. Cada etiqueta recibe una confianza según
su evidencia:

\begin{center}
\small
\begin{tabular}{lc@{\hspace{2.5em}}lc}
  \toprule
  Evidencia & Confianza & Evidencia & Confianza \\
  \midrule
  Positivo en la conclusión & 0,95 & Mención incierta & 0,50 \\
  Positivo en el cuerpo & 0,85 & Negado & 0,95 \\
  Evidencia débil & 0,70 & Sin mención & 0,80 \\
  \bottomrule
\end{tabular}
\end{center}

Las reglas se construyeron y auditaron solo con reportes no dorados, y se congelaron (etiqueta
de git) antes de compararlas con los dorados.

\textbf{Validación del extractor (entregable 1).} Con las menciones inciertas como positivas
(U-ones), la exactitud por pares estudio$\times$etiqueta es 0,810 [0,78--0,84], lo que
equivale a una \textbf{tasa de ruido del 19\,\%}. El F1 macro es 0,72 [0,66--0,76], el $\kappa$
macro 0,56 [0,49--0,62] y la AUC macro del objetivo blando 0,81 [0,78--0,85]; los IC 95\,\%
son por bootstrap de estudios.

Las menciones inciertas resultaron correctas solo el 36\,\% de las veces, así que para
entrenar usamos U-zeros (exactitud 0,822). La confianza está sobrestimada: los positivos en la
conclusión aciertan el 76\,\% de las veces, frente al 0,95 asignado.

La Tabla~\ref{tab:etiquetas} muestra que el texto es fiable para el LCA y los meniscos, pero
no para el derrame ni la sinovitis ($\kappa \approx 0{,}3$). Los dorados parecen anotados
sobre la imagen, y esos hallazgos se omiten a menudo en el reporte.

\textbf{Desbalance (entregable 6).} La Tabla~\ref{tab:etiquetas} cuantifica la prevalencia
de cada hallazgo en la muestra. Fractura, sinovitis, contusión y menisco lateral están por
debajo del 15\,\%. La mitigación usa \texttt{pos\_weight}$_l = N^-_l/N^+_l$, que va de 0,9
a 9,2.

Los dorados están enriquecidos respecto de la muestra (sinovitis 47\,\% frente a 10\,\%).
Esto no afecta a la AUC, pero sí a la precisión y al recall de cualquier umbral fijado sobre
datos débiles.

\begin{table}[t]
  \caption{Etiquetas por hallazgo. Prevalencia de la etiqueta débil U-zeros en
    \texttt{weak\_train} ($n=510$) y de la etiqueta dorada en \texttt{gold\_val} ($n=58$);
    peso de positivos usado en la pérdida; acuerdo del extractor con los dorados (F1 y
    $\kappa$, U-zeros).}
  \label{tab:etiquetas}
  \centering
  \small
  \begin{tabular}{lrrrcc}
    \toprule
    Hallazgo & Prev. débil (\%) & Prev. dorada (\%) & \texttt{pos\_weight} & F1 texto & $\kappa$ texto \\
    \midrule
    LCA               & 17,6 & 41,4 & 4,67 & 0,88 & 0,79 \\
    LCM               & 15,5 & 15,5 & 5,46 & 0,64 & 0,55 \\
    Menisco medial    & 38,6 & 44,8 & 1,59 & 0,84 & 0,72 \\
    Menisco lateral   & 14,1 & 39,7 & 6,08 & 0,81 & 0,70 \\
    Artrosis medial   & 22,7 & 25,9 & 3,40 & 0,69 & 0,61 \\
    Artrosis lateral  & 15,9 & 19,0 & 5,30 & 0,60 & 0,52 \\
    Artrosis PF       & 22,5 & 36,2 & 3,43 & 0,65 & 0,54 \\
    Derrame           & 52,7 & 60,3 & 0,90 & 0,79 & 0,33 \\
    Sinovitis         & 10,2 & 46,6 & 8,81 & 0,56 & 0,32 \\
    Quiste de Baker   & 25,9 & 20,7 & 2,86 & 0,76 & 0,68 \\
    Contusión         & 11,2 & 32,8 & 7,95 & 0,67 & 0,52 \\
    Fractura          &  9,8 & 31,0 & 9,20 & 0,70 & 0,56 \\
    \bottomrule
  \end{tabular}
\end{table}
